
\documentclass{article}
\usepackage{colortbl}
\usepackage{makecell}
\usepackage{multirow}
\usepackage{supertabular}

\begin{document}

\newcounter{utterance}

\centering \large Interaction Transcript for game `codenames', experiment `unambiguous', episode 2 with reflexion{-}llama{-}3.1{-}8b{-}instruct.
\vspace{24pt}

{ \footnotesize  \setcounter{utterance}{1}
\setlength{\tabcolsep}{0pt}
\begin{supertabular}{c@{$\;$}|p{.15\linewidth}@{}p{.15\linewidth}p{.15\linewidth}p{.15\linewidth}p{.15\linewidth}p{.15\linewidth}}
    \# & $\;$A & \multicolumn{4}{c}{Game Master} & $\;\:$B\\
    \hline

    \theutterance \stepcounter{utterance}  
    & & & \multicolumn{2}{p{0.3\linewidth}}{
        \cellcolor[rgb]{0.9,0.9,0.9}{
            \makecell[{{p{\linewidth}}}]{
                \texttt{\tiny{[GM$|$GM]}}
                {'hidden': {'team': ['coleslaw'], 'innocent': ['vinegar', 'cocaine', 'equality', 'propeller', 'bikini'], 'opponent': ['commie'], 'assassin': ['census']}, 'revealed': {'team': {'team': [], 'innocent': [], 'opponent': [], 'assassin': []}, 'opponent': {'team': [], 'innocent': [], 'opponent': [], 'assassin': []}}}
            }
        }
    }
    & & \\ \\

    \theutterance \stepcounter{utterance}  
    & & \multicolumn{4}{p{0.6\linewidth}}{
        \cellcolor[rgb]{0.9,0.9,0.9}{
            \makecell[{{p{\linewidth}}}]{
                \texttt{\tiny{[P1$\langle$GM]}}
                \texttt{You have attempted this same task before and did not succeed. These are the plans you wrote after those attempts.} \\
\texttt{Trial \#1: New plan: In my last attempt, I provided a clue that was too specific and related to only one word, coleslaw. This made it difficult for my teammate to guess the other remaining team words. Next time, I will provide a clue that is more general and relates to multiple words. I will also make sure to follow the format and rules of the game, avoiding any words that are not allowed or that do not fit the specified format. For example, I will not provide a clue that starts with "CLUE: condiment" and then list only one target word, as this is not a valid response. Instead, I will provide a clue that is more general and relates to multiple words, such as "CLUE: food" and then list the target words that are related to food, such as coleslaw, seafood, and peninsula.} \\
\\ 
\texttt{Trial \#2: New plan: In my last attempt, I provided a clue that was too specific and related to only one word, coleslaw. This made it difficult for my teammate to guess the other remaining team words. What's worse, I provided a clue that was not semantically related to the target words, as "salad" is not a general enough term to encompass all the words on the list. Next time, I will provide a clue that is more general and relates to multiple words. I will also make sure to follow the format and rules of the game, avoiding any words that are not allowed or that do not fit the specified format. For example, I will not provide a clue that starts with "CLUE: condiment" and then list only one target word, as this is not a valid response. Instead, I will provide a clue that is more general and relates to multiple words, such as "CLUE: food" and then list the target words that are related to food, such as coleslaw, seafood, and peninsula.} \\
\\ 
\texttt{Trial \#3: New plan: In my last attempt, I provided a clue that was too specific and related to only one word, coleslaw. This made it difficult for my teammate to guess the other remaining team words. What's worse, I provided a clue that was not semantically related to the target words, as "salad" is not a general enough term to encompass all the words on the list. I also failed to follow the format and rules of the game by including a distractor word, "peninsula", in the list of targets. Next time, I will provide a clue that is more general and relates to multiple words. I will also make sure to follow the format and rules of the game, avoiding any words that are not allowed or that do not fit the specified format. For example, I will provide a clue that is more general and relates to multiple words, such as "CLUE: food" and then list the target words that are related to food, such as coleslaw, seafood, and other food{-}related words. I will also make sure to only include team words in the list of targets, without any distractors.} \\
\\ 
\\ 
\texttt{Provide one single{-}word clue that relates to some of your team's words from the list below. You can choose to how many and to which words your clue relates to.} \\
\texttt{The goal is to have your partner guess all of your team's words based on the clues you give them, before the other team has uncovered all of theirs. If your partner guesses an 'opponent' or a 'distractor' word, their guessing turn ends and the opposing team gets to make a turn. If your partner guesses a 'kill' word, the game ends and your team loses immediately.} \\
\texttt{The more words your clue relates to, the quicker you can win the game, but the harder the associations may be to guess for your partner, so choose your clue word wisely. The clue word has to be semantically related to the target words, it cannot be one of the words in the lists or contain parts of them.} \\
\\ 
\texttt{Always give your single{-}word clue and your comma{-}separated list of related target words in the following format and make your answers as short as possible, never include any other text than is required in this form:} \\
\\ 
\texttt{CLUE: <WORD>} \\
\texttt{TARGETS: <TARGETS>} \\
\\ 
\texttt{Your team words are: coleslaw.} \\
\texttt{Your opponent's team words are: commie.} \\
\texttt{Distractor words are: vinegar, cocaine, equality, propeller, bikini.} \\
\texttt{The kill words are: census.} \\
            }
        }
    }
    & \\ \\

    \theutterance \stepcounter{utterance}  
    & \multicolumn{4}{p{0.6\linewidth}}{
        \cellcolor[rgb]{0.9,0.9,0.9}{
            \makecell[{{p{\linewidth}}}]{
                \texttt{\tiny{[P1$\rangle$GM]}}
                \texttt{CLUE: salad} \\
\texttt{TARGETS: coleslaw} \\
            }
        }
    }
    & & \\ \\

    \theutterance \stepcounter{utterance}  
    & & & \multicolumn{2}{p{0.3\linewidth}}{
        \cellcolor[rgb]{0.9,0.9,0.9}{
            \makecell[{{p{\linewidth}}}]{
                \texttt{\tiny{[GM$|$GM]}}
                \texttt{salad} \\
            }
        }
    }
    & & \\ \\

    \theutterance \stepcounter{utterance}  
    & & & \multicolumn{2}{p{0.3\linewidth}}{
        \cellcolor[rgb]{0.9,0.9,0.9}{
            \makecell[{{p{\linewidth}}}]{
                \texttt{\tiny{[GM$|$GM]}}
                ['coleslaw']
            }
        }
    }
    & & \\ \\

    \theutterance \stepcounter{utterance}  
    & & & \multicolumn{2}{p{0.3\linewidth}}{
        \cellcolor[rgb]{0.9,0.9,0.9}{
            \makecell[{{p{\linewidth}}}]{
                \texttt{\tiny{[GM$|$GM]}}
                {'word': 'coleslaw', 'assignment': 'team'}
            }
        }
    }
    & & \\ \\

    \theutterance \stepcounter{utterance}  
    & & \multicolumn{4}{p{0.6\linewidth}}{
        \cellcolor[rgb]{0.9,0.9,0.9}{
            \makecell[{{p{\linewidth}}}]{
                \texttt{\tiny{[GM$\rangle$P2]}}
                \texttt{You have attempted this same task before and did not succeed. These are the plans you wrote after those attempts.} \\
\texttt{Trial \#1: New plan: In my last attempt, I made a mistake by including an opponent word ("stead") in my guess. Next time, I will make sure to exclude opponent words from my list. I will also pay closer attention to the list of words provided and only choose words that are closely associated with the target word. For example, if the target word is "condiment", I will choose words like "vinegar" or "coleslaw" that are directly related to condiments.} \\
\\ 
\texttt{Trial \#2: New plan: In my last attempt, I made a mistake by including two words in my guess instead of a single list of words. Next time, I will make sure to follow the format exactly and provide a single list of words that best relate to the target word "salad". I will also pay closer attention to the list of words provided and only choose words that are closely associated with the target word. For example, if the target word is "salad", I will choose words like "coleslaw" that are directly related to salads.} \\
\\ 
\\ 
\texttt{Provide a comma{-}separated list of up to 1 words from the following list that best relate or are most closely associated with the word 'salad'. Always start your list of guess(es) with 'GUESS: ' and do not include any other text in your answer.} \\
\\ 
\texttt{census, bikini, vinegar, commie, propeller, coleslaw, cocaine, equality} \\
            }
        }
    }
    & \\ \\

    \theutterance \stepcounter{utterance}  
    & & & \multicolumn{4}{p{0.6\linewidth}}{
        \cellcolor[rgb]{0.9,0.9,0.9}{
            \makecell[{{p{\linewidth}}}]{
                \texttt{\tiny{[GM$\langle$P2]}}
                \texttt{GUESS: coleslaw} \\
            }
        }
    }
     \\ \\

    \theutterance \stepcounter{utterance}  
    & & & \multicolumn{2}{p{0.3\linewidth}}{
        \cellcolor[rgb]{0.9,0.9,0.9}{
            \makecell[{{p{\linewidth}}}]{
                \texttt{\tiny{[GM$|$GM]}}
                ['coleslaw']
            }
        }
    }
    & & \\ \\

    \theutterance \stepcounter{utterance}  
    & & & \multicolumn{2}{p{0.3\linewidth}}{
        \cellcolor[rgb]{0.9,0.9,0.9}{
            \makecell[{{p{\linewidth}}}]{
                \texttt{\tiny{[GM$|$GM]}}
                {'word': 'coleslaw', 'assignment': 'team'}
            }
        }
    }
    & & \\ \\

    \theutterance \stepcounter{utterance}  
    & & & \multicolumn{2}{p{0.3\linewidth}}{
        \cellcolor[rgb]{0.9,0.9,0.9}{
            \makecell[{{p{\linewidth}}}]{
                \texttt{\tiny{[GM$|$GM]}}
                {'word': 'coleslaw', 'assignment': 'team'}
            }
        }
    }
    & & \\ \\

    \theutterance \stepcounter{utterance}  
    & & & \multicolumn{2}{p{0.3\linewidth}}{
        \cellcolor[rgb]{0.9,0.9,0.9}{
            \makecell[{{p{\linewidth}}}]{
                \texttt{\tiny{[GM$|$GM]}}
                \texttt{team has won} \\
            }
        }
    }
    & & \\ \\

\end{supertabular}
}

\end{document}
